\chapter{Neue Drucksachen}
\quellseite{10}

Elfter Synodalbericht des California- und Nevada-Distrikts der deutschen ev.-luth. Synode von Missouri, Ohio u. a. Staaten. 1903. St. Louis, Mo. Concordia Publishing House. 62 Seiten 9×6. Preis: 15 Cts.

Gegenstand der Lehrverhandlungen war: „Der rechte Gebrauch der Gnadenmittel.“ Es wird zunächst gezeigt, was „Gnadenmittel“ seien, nämlich die von Gott selbst geordneten Mittel, wodurch die von Christo erworbene Gnade dem Sünder anbietet, im Glauben zueignet und versiegelt. Solche Mittel sind das Wort des Evangeliums und die Sacramente der Taufe und des Abendmahls. An diese Mittel hat Gott uns Menschen gebunden. So müssen denn die Menschen, sollen sie anders der von Christo erworbenen Gnade theilhaftig werden, zum Gebrauch der Gnadenmittel kommen. Gottes Gnade außerhalb der Gnadenmittel suchen ist Schwärmerei und führt jedes Mal zur Fälschung der Lehre von der Rechtfertigung, das heißt, der Gnadenstand wird, anstatt auf die von Christo erworbene und im Evangelium bezeugte gnädige Gesinnung Gottes, auf des Menschen Erneuerung und Heiligung gegründet. Hierauf wird der rechte Gebrauch der Gnadenmittel unter Berücksichtigung der Irrthümer der Secten und der Römischen sowie der Schwachheit der rechtgläubigen Christen allseitig auf Grund der Schrift dargelegt. Aus dem Bericht der Missionscommission geht hervor, daß die Sammelarbeit dieses Districts über ein weites Gebiet sich erstreckt, beschwerliche Arbeit, viel Geduld und bedeutende Geldauslagen erfordert, aber durch Gottes Gnade auch nicht vergeblich gewesen ist. Mehr Arbeiter werden dringend gewünscht. F. B.
